% Intended LaTeX compiler: pdflatex
\documentclass[a4paper,12pt]{article}
\usepackage[utf8]{inputenc}
\usepackage[T1]{fontenc}
\usepackage{graphicx}
\usepackage{longtable}
\usepackage{wrapfig}
\usepackage{rotating}
\usepackage[normalem]{ulem}
\usepackage{amsmath}
\usepackage{amssymb}
\usepackage{capt-of}
\usepackage{hyperref}
\usepackage[margin=0.2in]{geometry}
\usepackage{amsmath}
\usepackage{amsfonts}
\usepackage{indentfirst}
\usepackage{pgffor}
\usepackage{amssymb}
\usepackage{pifont}
\usepackage{cancel}
\usepackage{amsthm}
\usepackage{polynom}
\usepackage{tikz}
\usepackage[tikz]{bclogo}
\usepackage{listings}
\usepackage[most]{tcolorbox}
\usepackage{forest}
\usepackage{adjustbox}
\usepackage{tikz-3dplot}
\usepackage{mathtools}
\usepackage{pgfplots}
\usetikzlibrary{tikzmark,calc,fit,matrix,arrows,automata,positioning}
\usepackage{centernot}
\usepackage{lmodern}
\usepackage{physics}
\usepackage[boxed,linesnumbered,vlined]{algorithm2e}
\usepackage{algpseudocode}
\usepackage{algorithmicx}
\usepackage{svg}
\tikzstyle{every picture}+=[remember picture]
\DeclareMathOperator{\spn}{span}
\DeclareMathOperator{\dom}{domain}
\DeclareMathOperator{\ran}{range}
\DeclareMathOperator{\rng}{range}
\DeclareMathOperator{\img}{Im}
\DeclareMathOperator{\adj}{adj}
\DeclareMathOperator{\Var}{Var}
\DeclareMathOperator{\cov}{cov}
\newcommand\dif[1]{\:\textrm{d}#1}
\DeclarePairedDelimiter\ceil{\lceil}{\rceil}
\DeclarePairedDelimiter\floor{\lfloor}{\rfloor}
\newcommand{\ihat}{\hat{\textbf{\i}}}
\newcommand{\jhat}{\hat{\textbf{\j}}}
\newcommand{\khat}{\hat{\textbf{k}}}
\newcommand{\rhat}{\hat{\textbf{r}}}
\newcommand{\thetahat}{\boldsymbol{\hat{\theta}}}
\newcommand{\?}{\stackrel{?}{=}}
\newcommand{\cmark}{\text{\ding{51}}}
\newcommand{\xmark}{\text{\ding{55}}}
\author{mahmood}
\date{\today}
\title{im crazy}
\hypersetup{
 pdfauthor={mahmood},
 pdftitle={im crazy},
 pdfkeywords={},
 pdfsubject={im crazy},
 pdfcreator={Emacs 30.0.50 (Org mode 9.6)}, 
 pdflang={English}}
\begin{document}

\maketitle
\definecolor{Brown}{cmyk}{0,0.81,1,0.60}
\definecolor{OliveGreen}{cmyk}{0.64,0,0.95,0.40}
\definecolor{CadetBlue}{cmyk}{0.62,0.57,0.23,0}
\definecolor{lightlightgray}{gray}{0.9}
\lstset{
  language=C,                             % Code langugage
  basicstyle=\ttfamily,                   % Code font, Examples: \footnotesize, \ttfamily
  keywordstyle=\color{OliveGreen},        % Keywords font ('*' = uppercase)
  commentstyle=\color{gray},              % Comments font
  numbers=left,                           % Line nums position
  numberstyle=\tiny,                      % Line-numbers fonts
  stepnumber=1,                           % Step between two line-numbers
  numbersep=5pt,                          % How far are line-numbers from code
  backgroundcolor=\color{lightlightgray}, % Choose background color
  frame=single,                           % A frame around the code
  tabsize=2,                              % Default tab size
  captionpos=b,                           % Caption-position = bottom
  breaklines=true,                        % Automatic line breaking?
  breakatwhitespace=false,                % Automatic breaks only at whitespace?
  showspaces=false,                       % Dont make spaces visible
  showtabs=false,                         % Dont make tabls visible
  columns=flexible,                       % Column format
  morekeywords={__global__, __device__},  % CUDA specific keywords
  showstringspaces=false,                 % dont show spaces in strings
}
% so that latex doesnt complain about sage not being a language
\lstdefinelanguage{sage}{}

% environment definition for special blocks
\theoremstyle{definition}
\newtheorem{my_example}{Example}
\counterwithin{my_example}{section}
\counterwithin{my_example}{subsection}
\newtheorem{definition}{Definition}
\counterwithin{definition}{section}
\counterwithin{definition}{subsection}
\newtheorem{characteristic}{Characteristic}
\newtheorem{entailment}{Entailment}
%\newtheore*{proof}{Proof}
\newtheorem{lemma}{Lemma}
\newtheorem{question}{Question}
\newtheorem{subquestion}{Subquestion}[question]
%\newtheorem{note}{Note}
\newtheorem{answer}{Answer}
\counterwithin{answer}{question}
\counterwithin{answer}{subquestion}
\newtheorem{step}{Step}[answer]

% the following snippet auto resizes images to fit properly
% Determine if the image is too wide for the page.
% \makeatletter
% \def\ScaleIfNeeded{%
%   \ifdim\Gin@nat@width>\linewidth
%     \linewidth
%   \else
%     \Gin@nat@width
%   \fi
% }
% \makeatother
% % Resize figures that are too wide for the page.
% \let\oldincludegraphics\includegraphics
% \renewcommand\includegraphics[2][]{%
%   \oldincludegraphics[width=\ScaleIfNeeded]{#2}
% }

\newenvironment{note}
  {\begin{bclogo}[logo=\bcinfo, couleurBarre=orange, couleur=lightgray]{ note}}
  {\end{bclogo}}
\newenvironment{attention}
  {\begin{bclogo}[logo=\bcattention, couleurBarre=red, noborder=true, 
                couleur=red!15]{Important!}}
  {\end{bclogo}}

so i was a bit curious about the size of the files ive been manually writing for the past year since i started college\\[0pt]
and boy was i shocked to see the outcome of this curiousity\\[0pt]
\section{number of files}
\label{sec:org755c064}
lets start with the least interesting part, the number of files ive created and edited and still have.\\[0pt]
i write almost everything in \texttt{tex} or \texttt{org} files so thats what im gonna search for here.\\[0pt]
\begin{lstlisting}[language=bash,label= ,caption= ,captionpos=b,numbers=none]
echo in total, we have: $(find ../ -name '*.org' -or -name '*.tex' | wc -l) files
\end{lstlisting}

\begin{verbatim}
in total, we have: 709 files
\end{verbatim}
\section{actual file list}
\label{sec:org12974fe}
now we list them and check how much text each file contains:\\[0pt]
\begin{lstlisting}[language=bash,label= ,caption= ,captionpos=b,numbers=none]
find .. -name '*.org' -or -name '*.tex' -not -path '*logseq*' -not -path '*out/*' | tr '\n' '\0' | xargs -0 wc | sort -k3 -n -r | head -30
\end{lstlisting}

\begin{center}
\begin{tabular}{rrrl}
lines & words & characters & filename\\[0pt]
\hline
65961 & 303085 & 2279876 & total\\[0pt]
3584 & 15739 & 143447 & ../data\textsubscript{structures}/data\textsubscript{structures.org}\\[0pt]
2585 & 16032 & 116349 & ../linear\textsubscript{algebra2}/linear\textsubscript{algebra2.org}\\[0pt]
2296 & 11418 & 96609 & ../calculus2/calculus2.org\\[0pt]
1239 & 4764 & 59047 & ../data\textsubscript{structures}/homework8.org\\[0pt]
945 & 5966 & 48088 & ../networking/networking.org\\[0pt]
953 & 4947 & 37896 & ../discrete\textsubscript{maths2}/discrete\textsubscript{maths2.org}\\[0pt]
1197 & 5363 & 34624 & ../linear\textsubscript{algebra2}/homework5.org\\[0pt]
1093 & 5444 & 33665 & ../notes/data/FC/03889B-04C0-453D-B178-987A12945095/boolean\textsubscript{algebra.tex}\\[0pt]
895 & 3044 & 32651 & ../code/tikz.org\\[0pt]
877 & 5060 & 31511 & ../notes/boolean\textsubscript{algebra.org}\\[0pt]
280 & 1819 & 30552 & ../notes/wallpaper.org\\[0pt]
1169 & 4877 & 29905 & ../linear\textsubscript{algebra2}/homework9.org\\[0pt]
155 & 2110 & 29708 & ../notes/movies.org\\[0pt]
479 & 3366 & 25742 & ../notes/agenda.org\\[0pt]
1053 & 5916 & 25682 & ../discrete\textsubscript{maths}/homework11.tex\\[0pt]
818 & 3084 & 25209 & ../code/sage.org\\[0pt]
1229 & 6328 & 23173 & ../linear\textsubscript{algebra}/homework4.tex\\[0pt]
745 & 2491 & 22410 & ../gui\textsubscript{programming}/gui\textsubscript{programming.org}\\[0pt]
1158 & 3431 & 21523 & ../linear\textsubscript{algebra}/homework11.tex\\[0pt]
1929 & 3640 & 21096 & ../linear\textsubscript{algebra}/homework9.tex\\[0pt]
337 & 2519 & 18441 & ../linear\textsubscript{algebra2}/homework6.org\\[0pt]
254 & 2514 & 17158 & ../linear\textsubscript{algebra2}/homework8.org\\[0pt]
431 & 3660 & 16063 & ../discrete\textsubscript{maths}/homework10.tex\\[0pt]
335 & 845 & 15976 & ../data\textsubscript{structures}/homework6.org\\[0pt]
336 & 1924 & 15587 & ../notes/20230101161848-algorithms\textsubscript{homework}\textsubscript{6.org}\\[0pt]
244 & 1699 & 15427 & ../notes/algorithms.org\\[0pt]
596 & 2227 & 15019 & ../linear\textsubscript{algebra2}/homework3.tex\\[0pt]
343 & 2015 & 14924 & ../notes/20221204105120-combinatorics.org\\[0pt]
371 & 1855 & 14784 & ../notes/20221105001640-recursive\textsubscript{function.org}\\[0pt]
\end{tabular}
\end{center}

thats right, more than \textbf{2 million} characters ive written \textbf{and} saved , who knows how much i have actually written, then edited or deleted, but this is what i have so far.\\[0pt]
\begin{note}
i did trim the file list because theres some private stuff in there\\[0pt]
\end{note}
\end{document}